\subsection{Đo lường độ phù hợp của mô hình}

Việc đánh giá một mô hình hồi quy bắt đầu từ khả năng "khớp" của nó đối với dữ liệu thực tế. Để thực hiện điều này, các nhà thống kê sử dụng phương pháp phân rã tổng phương sai để hiểu rõ nguồn gốc của sự biến động trong biến mục tiêu.

Để hiểu về $R^2$, trước hết ta cần phân rã sự biến thiên của biến phụ thuộc. Tổng bình phương toàn phần (TSS) được chia thành hai thành phần:
\begin{equation}
    \underbrace{\sum (y_i - \bar{y})^2}_{TSS} = \underbrace{\sum (\hat{y}_i - \bar{y})^2}_{ESS} + \underbrace{\sum (y_i - \hat{y}_i)^2}_{RSS}
\end{equation}
Trong đó:
\begin{itemize}
    \item \textbf{TSS (Total Sum of Squares):} Tổng bình phương toàn phần, đo lường sự biến thiên tổng thể của biến phụ thuộc xung quanh giá trị trung bình mẫu. Đây chính là mức độ rủi ro hoặc sự không chắc chắn ban đầu khi chưa có sự hỗ trợ của các biến độc lập.
    \item \textbf{ESS (Explained Sum of Squares):} Tổng bình phương được giải thích, đại diện cho phần biến thiên mà mô hình đã nắm bắt được thông qua các đặc trưng đầu vào. Một mô hình lý tưởng là mô hình có ESS chiếm tỷ trọng lớn trong TSS.
    \item \textbf{RSS (Residual Sum of Squares):} Tổng bình phương phần dư, là phần sai số còn sót lại mà mô hình chưa thể giải thích. Phần dư này thường được kỳ vọng là những nhiễu trắng ngẫu nhiên không mang thông tin hệ thống.
\end{itemize}

\subsubsection{Hệ số xác định ($R^2$)}
$R^2$ đại diện cho tỷ lệ phần trăm sự biến động của biến phụ thuộc được giải thích bởi các biến độc lập:
\begin{equation}
R^2 = \frac{ESS}{TSS} = 1 - \frac{RSS}{TSS}
\end{equation}

Về mặt trực giác, $R^2$ cung cấp một thang đo không đơn vị, cho phép chúng ta so sánh mức độ giải thích của các mô hình khác nhau trên cùng một tập dữ liệu. Một giá trị $R^2$ gần bằng 1 cho thấy mô hình đã giải thích được hầu hết sự biến động, trong khi giá trị gần bằng 0 cho thấy các biến độc lập cung cấp rất ít thông tin về biến phụ thuộc.

\textbf{Hạn chế:} Tuy nhiên, việc lạm dụng $R^2$ như một tiêu chí duy nhất để chọn mô hình là một sai lầm phổ biến. Một vấn đề cố hữu của $R^2$ là nó mang tính "tham lam" về mặt toán học. Khi ta thêm bất kỳ biến nào vào mô hình (ngay cả những biến không liên quan như "lượng mưa" để dự báo "giá chứng khoán"), $R^2$ vẫn sẽ tăng hoặc giữ nguyên do RSS luôn giảm theo xu hướng đại số khi bậc tự do giảm xuống. Điều này dẫn đến nguy cơ quá khớp (\textit{overfitting}), khiến mô hình hoạt động rất tốt trên dữ liệu huấn luyện nhưng lại mất khả năng dự báo trên dữ liệu thực tế mới.

\subsubsection{Hệ số xác định hiệu chỉnh ($\bar{R}^2$)}
Để khắc phục nhược điểm trên, $\bar{R}^2$ đưa vào một "hình phạt" cho việc thêm các biến không thực sự có ý nghĩa bằng cách điều chỉnh theo bậc tự do:
\begin{equation}
\bar{R}^2 = 1 - \frac{RSS/(n - p - 1)}{TSS/(n - 1)}
\end{equation}
Khác với $R^2$ truyền thống, hệ số hiệu chỉnh này không nhất thiết phải tăng khi số lượng biến độc lập tăng lên. Nó đóng vai trò như một bộ lọc thông minh; giá trị này chỉ tăng khi biến mới cải thiện mô hình đủ nhiều để bù đắp cho sự sụt giảm của bậc tự do. Do đó, trong thực tế phân tích hồi quy đa biến, các chuyên gia thường ưu tiên sử dụng $\bar{R}^2$ để lựa chọn mô hình tối ưu, đảm bảo tính tiết kiệm (parsimony) mà vẫn giữ được sức mạnh giải thích.

\subsection{Các giả định của phương pháp Bình phương tối thiểu (OLS)}
Để các ước lượng từ phương pháp Bình phương tối thiểu có các thuộc tính tốt như không chệch và hiệu quả nhất (BLUE), mô hình cần thỏa mãn các giả định nghiêm ngặt của định lý Gauss-Markov. Việc vi phạm bất kỳ giả định nào trong số này đều có thể khiến kết quả kiểm định $t$ và $F$ không còn giá trị tin cậy.

\begin{enumerate}
    \item \textbf{Tính tuyến tính:} Giả định này phát biểu rằng mối quan hệ giữa $y$ và các $x$ là tuyến tính đối với các tham số $\beta$. Điều này không có nghĩa là đồ thị phải là đường thẳng, mà là các tham số phải ở dạng bậc nhất, cho phép ta sử dụng đại số tuyến tính để giải bài toán tối ưu.
    \item \textbf{Mẫu ngẫu nhiên:} Dữ liệu được thu thập ngẫu nhiên từ quần thể, đảm bảo rằng kết quả ước lượng từ mẫu có thể suy rộng ra cho tổng thể mà không bị sai lệch do chọn lọc mẫu không đại diện.
    \item \textbf{Không có đa cộng tuyến hoàn hảo:} Các biến độc lập không được có mối quan hệ tuyến tính hoàn hảo với nhau. Nếu giả định này bị vi phạm, ma trận $(X^T X)$ sẽ không thể nghịch đảo, dẫn đến việc không thể xác định duy nhất các hệ số hồi quy.
    \item \textbf{Trị trung bình sai số bằng 0:} $E(u|X) = 0$. Giả định này đảm bảo rằng không có biến số quan trọng nào bị bỏ sót mà lại có tương quan với các biến trong mô hình, giúp các ước lượng hệ số không bị chệch.
    \item \textbf{Phương sai sai số thuần nhất:} $Var(u|X) = \sigma^2$ (không đổi). Sai số phải có độ phân tán đồng đều tại mọi mức độ của biến độc lập. Nếu phương sai thay đổi (heteroscedasticity), các sai số chuẩn sẽ bị sai lệch, khiến các kiểm định giả thuyết trở nên không chính xác.
    \item \textbf{Sai số có phân phối chuẩn:} $u \sim \mathcal{N}(0, \sigma^2)$. Đây là điều kiện cần thiết để thực hiện các kiểm định mẫu nhỏ. Trong trường hợp mẫu đủ lớn, nhờ định lý giới hạn trung tâm, giả định này có thể được nới lỏng nhưng vẫn là tiêu chuẩn vàng cho sự chính xác.
\end{enumerate}

\subsection{Kiểm định giả thuyết và Khoảng tin cậy}

Sau khi xây dựng mô hình, bước tiếp theo là xác định xem các mối quan hệ quan sát được là do bản chất quy luật hay chỉ là sự trùng hợp ngẫu nhiên trong dữ liệu mẫu.

\subsubsection{Kiểm định $t$ cho từng hệ số}
Kiểm định này giúp xác định liệu một biến cụ thể có đóng góp ý nghĩa vào mô hình hay không khi có sự hiện diện của các biến khác.
\begin{itemize}
    \item $H_0: \beta_j = 0$ (Biến $x_j$ không có tác động thực sự đến $y$).
    \item $H_1: \beta_j \neq 0$ (Biến $x_j$ có tác động đáng kể thống kê).
\end{itemize}
Về mặt thực hành, giá trị thống kê $t$ đo lường khoảng cách từ hệ số ước lượng đến giá trị 0 theo đơn vị sai số chuẩn. Nếu khoảng cách này đủ lớn, ta có bằng chứng để bác bỏ giả thuyết không.

\textbf{Khoảng tin cậy:} Một cách tiếp cận khác trực quan hơn là xây dựng khoảng tin cậy $(1-\alpha)$ cho $\beta_j$:
\begin{equation}
\hat{\beta}_j \pm t_{\alpha/2, n-p-1} \cdot se(\hat{\beta}_j)
\end{equation}
Khoảng tin cậy này cung cấp một vùng các giá trị khả dĩ cho tham số thực của quần thể. Nếu khoảng này không chứa giá trị 0, ta có thể kết luận một cách tự tin rằng biến đó có ý nghĩa thống kê ở mức ý nghĩa tương ứng. Nó cũng cho thấy độ chính xác của ước lượng; khoảng càng hẹp thì ước lượng càng tin cậy.

\subsubsection{Kiểm định $F$ cho toàn bộ mô hình}
Trong khi kiểm định $t$ tập trung vào từng biến đơn lẻ, kiểm định $F$ lại mang tính bao quát hơn, giúp trả lời câu hỏi: "Liệu tập hợp các biến độc lập này có giải thích được bất kỳ phần nào của biến phụ thuộc hay không?".
\begin{equation}
F = \frac{R^2/p}{(1 - R^2)/(n - p - 1)}
\end{equation}
Kiểm định $F$ so sánh sức mạnh giải thích của mô hình hiện tại với một "mô hình rỗng" (mô hình chỉ có hệ số chặn). Nếu giá trị $P\text{-value}$ nhỏ hơn mức ý nghĩa $\alpha$ (thường là 0.05), ta bác bỏ $H_0$ và khẳng định rằng ít nhất một trong các biến độc lập có tác động đáng kể. Điều này đặc biệt quan trọng khi các biến độc lập có tương quan cao với nhau, khiến từng biến riêng lẻ có thể không có ý nghĩa ($t$-test yếu) nhưng cả nhóm lại có sức mạnh giải thích rất lớn ($F$-test mạnh).

\subsection{Kết luận và So sánh}
Việc đánh giá mô hình hồi quy là một quá trình đa chiều, đòi hỏi sự kết hợp hài hòa giữa các chỉ số định lượng và việc kiểm tra các giả định lý thuyết. Một sai lầm phổ biến là quá tập trung vào một chỉ số duy nhất mà bỏ qua bức tranh tổng thể. Chẳng hạn, một mô hình có $R^2$ rất cao nhưng vi phạm giả định về phương sai sai số sẽ dẫn đến các khoảng tin cậy vô nghĩa. Ngược lại, một biến có hệ số nhỏ nhưng có ý nghĩa thống kê cao vẫn có thể mang lại những hiểu biết quý giá về mặt kinh tế hoặc xã hội.

Tóm lại, quá trình thẩm định mô hình cần tuân thủ một lộ trình nghiêm ngặt: đầu tiên là kiểm tra các giả định Gauss-Markov để đảm bảo tính hợp lệ, sau đó đánh giá ý nghĩa tổng thể qua kiểm định $F$, và cuối cùng là tinh chỉnh các biến thành phần qua kiểm định $t$ và hệ số $\bar{R}^2$.

\vspace{0.5cm}
\begin{table}[h]
\centering
\caption{Tóm tắt các chỉ số đánh giá và ý nghĩa ứng dụng}
\begin{tabular}{l p{6cm} l}
\toprule
\textbf{Chỉ số} & \textbf{Công dụng chính trong thực tiễn} & \textbf{Ngưỡng kỳ vọng} \\ \midrule
$R^2$ & Đo lường tỷ lệ biến thiên được giải thích. & Càng gần 1 càng tốt. \\
$\bar{R}^2$ & Lựa chọn số lượng biến tối ưu cho mô hình. & Tối đa hóa giá trị này. \\
$t$-test & Xác định vai trò riêng biệt của từng biến. & $|t| >$ giá trị tới hạn. \\
$F$-test & Xác nhận mô hình có giá trị sử dụng hay không. & $p$-value < 0.05. \\
$se(\hat{\beta})$ & Đánh giá độ ổn định của các hệ số ước lượng. & Càng nhỏ càng tốt. \\ \bottomrule
\end{tabular}
\end{table}
